% Transcription — fonds Alexandre Grothendieck, shelfmark 126, batch 1 (pages 1-20).
% Established from the facsimile archives/batches/126/batch-01.pdf (prepared
% from a copy of 126.pdf the user uploaded; PDF page k = archivists' page k,
% no cover sheet), per the skill transcribe-grothendieck. The folder is
% thirty-four pages; this is the first of two batches (pages 21-34 are batch
% 2, transcribed separately and in parallel by another pass).
%
% Dating: the folder is « s.d. » in src/content/catalogue.ts; its inventory
% group « Descentes » (126-128, cotes 126 to 133) is dated [avant 1970], and
% \dating says both, per the skill's group rule. Nothing on these leaves
% dates them. Pages 17-20 are written on sheets bearing a printed letterhead
% of a typing business; it is not his and not recorded.
%
% Pages skipped: 16, a blank white sheet. Page 15 is a tan cover inscribed
% in ink, top right and underlined, « Descente fidèlement plate »; it takes
% no \page marker and its words become a section heading with a \note.
% Page 8 is a tan sheet carrying only three lines of plan (« Généralités sur
% la topologie fid. plate / Les H^p(S'/S, H^q(F)) pour F quasi-cohérent /
% Applications aux relations d'équivalence ... »); it is transcribed, since
% it is his and about the mathematics.
%
% Pages summarised: none. Pages transcribed from his hand: 1-14, 17-20.
% His pagination: none visible on these leaves.
%
% What the batch is about. Three runs.
% (1) Pages 1-7: Čech cohomology and faithfully flat descent of
% quasi-coherent sheaves. Th. 1: for f : X -> S faithfully flat
% quasi-compact (quasi-separated, S affine) and F quasi-coherent, the Čech
% cohomology of X./S with coefficients in F is H^0(S,F) in degree 0 and zero
% otherwise; Cor. 1 (sheafified form), Cor. 2 (S' |-> Γ(S', F ⊗ S') is a
% sheaf for the fpqc topology; written sideways in the margin), Cor. 3 (for
% any topology between Zariski and fpqc, H*(X^T, F^T) = H*(X, F)), Cor. 4
% (the Čech-to-derived spectral sequence Ȟ^p(S'/S, H^q(F)) => H*(S,F)),
% the low-degree exact sequence and remarks. Pages 4-7: an application to
% infinitesimal deformation of an equivalence relation: given R_0 ⇉ X_0
% effective with X_0 -> Y_0 = X_0/R_0 fpqc, and a nilpotent thickening
% R ⇉ X, show that A = Ker(H^0(X,O) ⇉ H^0(R,O)) surjects onto A_0 (the
% obstruction ξ lives in Ȟ^0(X_0/Y_0, H^1(J)) which vanishes by the above),
% and that X -> Y is flat with R = X ×_Y X.
% (2) Pages 9-14: topology generated by a pretopology T on C and a
% pretopology T' on a full subcategory C' (with conditions (a), (b), and a
% struck (c)); Prop. 1 describes the covering families of the generated
% topology and checks the axioms; Cor. 2 (sheaves for it = T-sheaves
% inducing T'-sheaves on C'); Prop. 2, the case (Sch) with Zariski on it and
% finite flat surjective families on affines — the resulting topology,
% covering criteria (fpqc morphisms, flat families of open morphisms), and
% the sheaf criterion a) Zariski sheaf, b) exactness of
% F(S) -> F(S') ⇉ F(S' ×_S S') for S' -> S fpqc between affines.
% (3) Pages 17-20, after the cover « Descente fidèlement plate »: a plan
% (A) application to comparison with holomorphic étale varieties, B) to a
% canonical factorisation), then a Proposition: f : X -> Y finite
% (universally open), Y locally noetherian, with constant separable
% geometric degree d(y), factors as a finite surjective radicial map
% followed by a finite étale one, universal for Y-maps to étale Y-schemes.
% Proof sketches: Y' = Hom(I, X) = X^I and the locus Y'' where
% I × Spec k(y') -> X'_{y'} is bijective (p. 18, cancelled by two diagonals,
% continued on p. 20); a section s of a finite X -> Y has s(Y) open and
% closed, giving X = X_(1) ⊔ X^(1) with degrees 1 and n-1 (p. 19).
%
% Legibility: pages 1, 9-13, 17-20 carry their statements and much of the
% prose; pages 2-7 are fast, dense drafting in which the formulas carry and
% most connective prose does not, and is left \ill{}.
%
% Words he underlines are set in \emph, as in batch 2. Notation fixed once
% for the batch: \check{H} for his Čech cohomology,
% \mathcal{H}^q(F) for the presheaf T |-> H^q(T, F ⊗ T) (he writes a script
% H), \underline{H} for the sheafified Čech groups of Cor. 1, \mathcal{J} for
% the ideal sheaves of the thickening, \mathcal{T} for the generated
% topology (his round T), (\mathrm{Sch}), (\mathrm{Ens}).
%
% Readings to check first: the hypotheses of Th. 1 and the boxed marginal
% beside it (p. 1); the sideways Cor. 2 (p. 1); almost all prose of pp. 2-7;
% the opening of p. 4 (hypotheses on R_0); the struck (c) on p. 9; the
% right-hand prose of p. 13; p. 17's « rang ... géométrique » and the
% exponent on X_0; p. 18's lemma statement and its last lines.
%
% Pass: Opus 5.5 (claude-opus-5-5), 2026-09-24 — first pass, unchecked against the pages by a human.

\documentclass[11pt,a4paper]{article}
% Le préambule commun à toutes les transcriptions du fonds Grothendieck.
%
% Il tient en une idée : **l'incertitude se compose, elle ne se résout pas.**
% Une transcription qui lisse un mot douteux a détruit la seule chose pour
% laquelle on la faisait — un lecteur doit pouvoir savoir, à chaque endroit,
% ce qui était sur le feuillet et ce qui a été ajouté. Les six macros qui
% suivent sont donc l'appareil critique, et non de la décoration.
%
% Les mêmes commandes sont reconnues par `scripts/render.mjs`, qui produit la
% vue de lecture du site. Une macro ajoutée ici doit l'être aussi là-bas,
% faute de quoi le rendu échouera bruyamment — ce qui est voulu : un rendu
% silencieusement faux est pire qu'un rendu absent.

\ProvidesPackage{grothendieck}[2026/08/08 Transcriptions du fonds Grothendieck]

\usepackage{amsmath,amssymb,amsthm}
\usepackage{mathrsfs}
\usepackage{stmaryrd}
% Les diagrammes commutatifs sont partout dans ce fonds : c'est la langue
% même de la géométrie algébrique de Grothendieck, et une transcription qui
% les rendrait en prose ne transcrirait rien.
\usepackage{tikz-cd}
% Français par défaut — c'est la langue des feuillets — mais l'anglais est
% chargé aussi : la traduction et le résumé sont des documents anglais, et la
% typographie française (espaces avant les deux-points) y serait une faute.
% Ces documents basculent par \selectlanguage{english} en tête de corps.
\usepackage[english,french]{babel}
\usepackage{fontspec}
\usepackage{geometry}
\geometry{a4paper,margin=25mm,marginparwidth=28mm}
\usepackage{marginnote}
\usepackage{xcolor}
% ulem, not soul. `soul`'s \st and \ul are built on a character-by-character
% scan that XeLaTeX's font machinery breaks: the document fails with an
% undefined control sequence rather than degrading. ulem does the same two jobs
% and is Unicode-engine safe. `normalem` stops it hijacking \emph, which the
% transcriptions use for Grothendieck's own emphasis.
\usepackage[normalem]{ulem}
\usepackage{hyperref}
\hypersetup{colorlinks=true,linkcolor=black,urlcolor=black,pdfborder={0 0 0}}

% --- Métadonnées du lot -----------------------------------------------------
% Reprises telles quelles par le script de rendu, qui les affiche en tête de
% la vue de lecture. Elles fixent ce que le fichier prétend couvrir : sans
% elles, une transcription flotte sans point d'attache dans un fonds de seize
% mille feuillets.

\newcommand{\folder}[1]{\gdef\grothFolder{#1}}
\newcommand{\batch}[1]{\gdef\grothBatch{#1}}
\newcommand{\leaves}[2]{\gdef\grothFirst{#1}\gdef\grothLast{#2}}
\newcommand{\foldertitle}[1]{\gdef\grothFoldertitle{#1}}
\gdef\grothFolder{?}\gdef\grothBatch{?}\gdef\grothFirst{?}\gdef\grothLast{?}\gdef\grothFoldertitle{}

% --- Le résumé --------------------------------------------------------------
%
% Il ouvre la lecture modernisée, sous le titre « Résumé », et il est composé
% en petit. Ce n'est pas un ornement : c'est le seul passage du document qui
% ne soit pas des mathématiques, et un lecteur doit pouvoir le reconnaître
% avant de l'avoir lu — puis le sauter s'il connaît déjà le sujet, ou s'y
% arrêter s'il ne le connaît pas. Le filet qui le suit marque la reprise du
% corps.

\newenvironment{resume}
  {\small\setlength{\parskip}{0.45em}}
  {\par\medskip\hrule\medskip}

% --- Mots-clés ---------------------------------------------------------------
%
% En anglais, en clôture du résumé de la lecture modernisée : le vocabulaire
% moderne sous lequel chercher ce que les feuillets construisent. Le manifeste
% du site (`npm run manifest`) les extrait pour étiqueter la cote entière —
% cette ligne est la source unique des tags, il n'y a pas de fichier à côté.
\newcommand{\keywords}[1]{\par\smallskip\noindent
  {\footnotesize\textsc{Keywords} --- \textit{#1}}\par}

% --- L'appareil critique ----------------------------------------------------

% Le numéro de feuillet, en marge. C'est l'ancre qui permet de régler un
% désaccord sur une formule en regardant *un* feuillet plutôt qu'une chemise
% de six cents. Le site s'en sert aussi pour tourner les pages du fac-similé
% quand on fait défiler la transcription.
\newcommand{\leaf}[1]{%
  \marginnote{\footnotesize\color{gray!70}\textsf{#1}}%
  \label{leaf:#1}\ignorespaces}

% Illisible : on ne devine pas. Un mot inventé qui se lit comme les autres est
% le pire résultat possible.
\newcommand{\ill}{\textcolor{red!60!black}{[\,\ldots\,]}}

% Lecture douteuse : le mot est donné, et signalé comme incertain.
\newcommand{\uncertain}[1]{\uline{#1}}

% Ajout de l'éditeur — une lettre manquante, un mot sous-entendu.
\newcommand{\add}[1]{\textcolor{blue!50!black}{[#1]}}

% Biffé par Grothendieck lui-même. Ce qu'il a rayé fait partie du document :
% une rature dit souvent où la pensée a changé de direction.
\newcommand{\struck}[1]{\sout{#1}}

% Note du transcripteur, sur l'état matériel du feuillet.
\newcommand{\note}[1]{\par\smallskip{\small\color{gray!60!black}\textit{[#1]}}\par\smallskip}

% Note marginale de Grothendieck, distincte du corps du texte.
\newcommand{\marginal}[1]{\par\smallskip\noindent\hspace{1em}%
  {\small\color{gray!60!black}#1}\par\smallskip}

% --- Titre ------------------------------------------------------------------

\AtBeginDocument{%
  \begin{center}
    {\large\bfseries Fonds Alexandre Grothendieck --- cote n\textsuperscript{o}~\grothFolder}\\[2pt]
    {\itshape\grothFoldertitle}\\[4pt]
    {\small feuillets \grothFirst--\grothLast\ (lot \grothBatch)}\\[2pt]
    {\footnotesize\color{gray!60!black}Transcription établie d'après le fac-similé numérisé
     par l'Université de Montpellier.\\
     Les lectures incertaines sont soulignées, les illisibles notées [\,\ldots\,],
     les ajouts entre crochets.}
  \end{center}
  \vspace{1em}}

\endinput


\folder{126}
\batch{1}
\pages{1}{20}
\dating{s.d. — le groupe « Descentes » (126 à 133) est daté [avant 1970]}
\watermark{Édition de démonstration}
\foldertitle{Descente fidèlement plate. Divers : notes manuscrites (s.d.)}

\begin{document}

\section*{Cohomologie de Čech et descente fidèlement plate}
\note{titre de l'éditeur, qui ne figure pas sur les feuillets}

\page{1}

\textbf{Théorème 1.} $f : X \to S$ un morphisme fid. plat quasi-compact
\add{(}\uncertain{quasi-séparé}, $S$ affine\add{)}, $F$ faisceau quasi-cohérent
sur \uncertain{$X$}. Alors
\[
  \check{H}^i(X^{\cdot}/S, F) \simeq
  \begin{cases}
    0 & \text{si } i \neq 0 \\
    H^0(S, F) & \text{si } i = 0 .
  \end{cases}
\]
\marginal{couvrant quasi-séparé / quasi-\struck{\ill{}} plat}\note{deux lignes
encadrées à droite de « Alors », qui semblent proposer de remplacer les
hypothèses « fid. plat quasi-compact » ; la parenthèse « quasi-séparé, $S$
affine » est ajoutée en interligne au-dessus de l'énoncé et reliée à lui par
un trait}

On peut l'énoncer sous la forme suivante :

\textbf{Corollaire 1.} Hypothèses précédentes, sans hyp. affine sur
\uncertain{$X$} cependant. Les $\underline{H}^i(X^{\cdot}/S, F)$
(faisceaux \add{qu. coh.} sur \uncertain{$X$}) sont nuls si $i \neq 0$, et
$F \to \underline{H}^0(X^{\cdot}/X, F)$ est un isomorphisme.
\note{« qu. coh. » en interligne ; le second argument est bien écrit $X^{\cdot}/X$}

Or la formation des $\underline{H}^i(X^{\cdot}/S, F)$ est compatible avec
extension plate de la base [où les $X^i/S$ sont quasi-compacts
quasi-séparés], et il suffit de prouver que les assertions dites sont vraies
après changement de base plat couvrant \add{les $S_\lambda$ affines}
$S_\lambda \to S$, \ill{} \ill{} \ill{} une famille qui soit telle que les
$X_\lambda/S_\lambda$ \uncertain{ont} des sections. D'où le résultat, car
pour $S$ affine et $X/S$ ayant une section, \add{c'est immédiat}.
\note{« les $S_\lambda$ affines » et « c'est immédiat » sont en interligne}

\marginal{\textbf{Corollaire 2.} Soit $X \to S$ \uncertain{couvrant}
\struck{\ill{}} ; $H^0(S, F) \simeq H^0(X/S, F)$. Plus généralement, le
foncteur \struck{\ill{}} $S' \mapsto \Gamma(S', F \otimes_S S')$ sur
$(\mathrm{Sch})^{\circ}_{/S}$ est un faisceau pour la topologie
\uncertain{fidèlement plate} \struck{\ill{}} \uncertain{quasi-compacte}.
[En effet, on est ramené à 2 vérifications, pour les deux types de familles
couvrantes\,: Zariski, et \ill{}]}\note{écrit en travers, le long de la
marge gauche, et encadré d'un trait ; la phrase entre crochets est entourée.
Dans l'angle, également en travers, « 2 \uncertain{compléter} \ill{} plus loin »}

\textbf{Corollaire 3.} Mettons sur $(\mathrm{Sch})$ une topologie $T$
intermédiaire entre celle de Zariski et la fpqc. Alors pour tout faisceau
quasi-cohérent $F$ sur $X$, $H^{*}(X^{T}_{\uncertain{\scriptstyle 1}}, F^{T})
= H^{*}(X, F)$.

\emph{Dém.} On est ramené \uncertain{par Leray}, pour
$(\mathrm{Sch})_{T} \to (\mathrm{Sch})_{\mathrm{Zar}}$, à prouver que pour
$X$ \struck{$X$} affine, on a

\page{2}

\[
  \mathcal{H}^i(X^{T}, F^{T}) = 0, \qquad
  \mathcal{H}^0(X^{T}, F^{T}) \simeq H^0(X, F).
\]
\struck{Pour} La \struck{\ill{}} relation \uncertain{seconde} résulte de la
définition de $F^{T}$. Pour la première, \struck{\ill{}} c'est \ill{} par
\ill{} \ill{}, \ill{} \ill{} \ill{} par suite spectrale \ill{} de Čech\,:
\ill{} \ill{} \ill{} pour tt \struck{recouvrement} \ill{} \ill{}
$\mathfrak{X} = (X_{\cdot} \to X)$, \ill{} \uncertain{fid. plats}\ill{},
$\check{H}^p(\mathfrak{X}, F) = 0$ pour tt $p \neq 0$.

Or \ill{} $F$ \ill{} \ill{} \ill{} \ill{}, \ill{} \ill{} \ill{}
$\mathfrak{X}$ \ill{} \ill{} \struck{\ill{}} \uncertain{fid. plat}
\struck{quasi-compact} \ill{} \ill{} \ill{} affine. \ill{} \ill{} \ill{}
\ill{} Th.~1.

Utilisons la \uncertain{suite spectrale de Leray}, \ill{}\,:

\textbf{Corollaire 4.} Soit $S' \to S$ couvrant (pour la \ill{}
\ill{}), $F$ \ill{} \ill{} sur $S$. Alors on a une suite spectrale
\[
  H^{*}(S, F) \Longleftarrow \check{H}^p(S'/S, \mathcal{H}^q(F))
\]
[où $\mathcal{H}^q(F)$ désigne le foncteur $T \mapsto H^q(T, F \otimes_S T)$
sur $(\mathrm{Sch})^{\circ}_{/S}$].

\page{3}

\note{en tête de page, au-dessus d'un trait\,: « l'isomorphisme
$H^0(S, F) \simeq H^0(S'/S, F)$ \uncertain{déjà connu}, et »}

Notons [la suite exacte en basses dimensions\,:
\[
  0 \to \check{H}^1(S'/S, F) \to H^1(S, F) \to
  \check{H}^0(S'/S, \mathcal{H}^1(F)) \to \check{H}^2(S'/S, F) \to
  H^2(S, F) .
\]

\textbf{Remarques} 1) Bien \uncertain{entendu}, on peut \ill{} \ill{} \ill{}
\ill{} \ill{} \ill{} (\ill{}).

2) La démonstration du Th.~1 \uncertain{montre} \ill{} \ill{} [\,\ill{} $S$
affine, \ill{}
\[
  \check{H}^p(S'/S, \mathcal{H}^q(F)) = 0 \quad \text{pour tt } p,
  \text{ si } q \neq 0 .
\]
\struck{\ill{} \ill{} \ill{} \ill{}} \ill{} \ill{}, on \ill{} \ill{}
\struck{\ill{}} !
\[
  \check{H}^0(S'/S, \mathcal{H}^q(F)) =
  \mathrm{Ker}\bigl(H^q(S', F) \rightrightarrows H^q(S' \times_S S', F)\bigr)
  = 0
\]
\struck{Je \ill{} \uncertain{résultats} \ill{} \ill{} \ill{}}
\note{deux traits horizontaux ferment la page}

\page{4}

\ill{} \ill{} \ill{} $X_0$. On suppose que le \ill{}
\struck{\uncertain{quotient} $X_0/R_0$} $R_0$ est (M) \uncertain{effective},
\ill{} \uncertain{morphisme}, \struck{\ill{} \ill{} \ill{}} \ill{} \ill{}
\ill{} \ill{}, et que \add{et que $p_1^0 : R_0 \to X_0$}
\struck{$R_0$ \ill{}} \uncertain{fid. plat} \uncertain{quasi-compact}, et
\uncertain{quasi-séparé}.
\note{l'ajout « et que $p_1^0 : R_0 \to X_0$ » est entouré et relié par un
trait ; au-dessus de la rature, un $R$ isolé}

Sous ces conditions, $R$ est \uncertain{M}-\struck{\ill{}} effective, et si
\struck{Et} si $X_0 \to Y_0$ est plat i.e. $R_0 \to X_0$ \ill{}

[On peut supposer $\mathcal{J}_{X_0}$ de carré nul.]\note{entouré}

\emph{Démonstration} (Soit $Y_0 = X_0/R_0$ dans \ill{}) $X_0 \to Y_0$ est
couvrant, \struck{\ill{}} \uncertain{quasi-compact} et
$R_0 \rightrightarrows X_0 \times_{Y_0} X_0$.

\struck{\ill{}} \ill{} : $\mathcal{J}_{X_0}$, $\mathcal{J}_{R_0}$, donc
$\mathcal{J}_{X_0}$ \ill{} \add{\uncertain{vient}} \struck{\ill{}} \ill{}
$R_0 \simeq X_0 \times_{Y_0} X_0$, [\,\ill{} \ill{} \ill{} \ill{}
$\mathcal{J}_{X_0}$ \ill{} \uncertain{propre} \struck{\ill{}}
$\mathcal{J}_{X_0}$\,], \ill{} \ill{} \ill{}, \ill{} \ill{} \ill{} \ill{}
\uncertain{effectivement}, \ill{} $X_0 \to Y_0$ \uncertain{couvrant} et
$R_0 \simeq X_0 \times_{Y_0} X_0$,

\page{5}

$\mathcal{J}_{Y_0}$ \uncertain{faisceau cohérent} sur $Y_0$.

\ill{} \uncertain{coefficients} de $R$, on \ill{}, \ill{} \ill{} \ill{}
affine. Soit donc \ill{} $Y_0$, et \struck{$A = \check{H}^0(X/\ill{},
\mathcal{O}_X)$}
\[
  A = \mathrm{Ker}\bigl(H^0(X, \mathcal{O}_X) \rightrightarrows
  H^0(R, \mathcal{O}_R)\bigr)
\]
\ill{} \ill{} $A \to A_0$, dans \ill{}
\[
  \mathrm{Ker}\bigl(H^0(X_0, \mathcal{J}_{X_0}) \rightrightarrows
  H^0(R_0, \mathcal{J}_{R_0})\bigr) = H^0(Y_0, \mathcal{J}_{Y_0}) =
  \mathcal{J}_{Y_0} .
\]
\struck{\ill{} \ill{} \ill{} \ill{}}
\[
  0 \to \mathcal{J}_{Y_0} \to A \to A_0 \to 0
\]
\ill{} \uncertain{exacte}, i.e. $A \to A_0$ \uncertain{surjectif}.

Pour le voir, \struck{prenons} soit $\varphi_0 \in A_0$, \ill{}
\add{\uncertain{son image}} dans $X_0$,
\struck{$f_0 \in H^0(X_0, \mathcal{O}_{X_0})$}
\[
  f_0 \in \mathrm{Ker}\bigl(H^0(X_0, \mathcal{O}_{X_0}) \rightrightarrows
  H^0(R_0, \mathcal{O}_{R_0})\bigr) ;
\]
il \uncertain{suffit} de \uncertain{voir} que $f_0$ se \uncertain{remonte}
en \ill{} $f \in H^0(X, \mathcal{O}_X)$. En effet, a priori il y a
\uncertain{une obstruction} à \uncertain{relever} \ill{}\,; \ill{} \ill{}
les deux images \ill{}\,; $\xi \in H^1(X_0, \mathcal{J}_{X_0})$\,;
\ill{}

\page{6}

\ill{} \ill{} \ill{} \ill{} \uncertain{des obstructions}, par $p_1^0$ et
$p_2^0$, \ill{} \ill{} \ill{} \ill{} les obstructions, \ill{} $=$
\uncertain{solution} $g_0 \in H^0(R_0, \mathcal{O}_{R_0})$ \ill{}
$(R, \mathcal{O}_R)$. D'où
\[
  \xi \in \mathrm{Ker}\bigl(H^1(X_0, \mathcal{J}_{X_0}) \rightrightarrows
  H^1(R_0, \mathcal{J}_{R_0})\bigr)
  = \check{H}^0\bigl(X_0/Y_0, \mathcal{H}^1(\mathcal{J}_{Y_0})\bigr)
\]
\uncertain{qui est nul}, \ill{} la 2\textsuperscript{ème} \ill{} \ill{}
\ill{}, $X_0 \to Y_0$ \uncertain{étant couvrant} \ill{} \uncertain{quasi-compact}
et \uncertain{quasi-séparé}. \ill{} $\xi = 0$, d'où l'existence \ill{} \ill{}
$f$.

\ill{} on \uncertain{trouve} \struck{$f \in \mathrm{Ker}$} $p_1(f) = p_2(f)$,
\ill{} \ill{} $f$ \ill{}, on $f$ \uncertain{provient} d'un $\varphi \in A$,
qui \struck{\ill{}} \add{\uncertain{relève}} \ill{} $\varphi_0$.
\note{« \uncertain{relève} » en interligne au-dessus d'un mot biffé}

\begin{tikzcd}[column sep=small]
A \arrow[r] \arrow[d] & H^0(X, \mathcal{O}_X) \arrow[r, bend left=8, "p_1"] \arrow[r, bend right=8] \arrow[d] & H^0(R, \mathcal{O}_R) \arrow[d] \\
A_0 \arrow[r] & H^0(X_0, \mathcal{O}_{X_0}) \arrow[r, bend left=8, "p_2"] \arrow[r, bend right=8] & H^0(R_0, \mathcal{O}_{R_0})
\end{tikzcd}

\note{les doubles flèches $\rightrightarrows$ du dessin sont rendues par deux
arcs\,; $p_1$ est écrit au-dessus de la première, $p_2$ entre les deux
lignes. Sous le diagramme, $\varphi_0$ sous $A_0$ et $f_0$ sous
$H^0(X_0, \mathcal{O}_{X_0})$ ; au-dessus, $f$ sur $H^0(X, \mathcal{O}_X)$}

À \uncertain{présent}, on \ill{} \ill{} \ill{} \ill{} donc que
$y\ (= p_2(f) - p_1(f)) \in H^0(R_0, \mathcal{J}_{R_0})$\,;
\uncertain{vérifions} que
\[
  y \in Z^1(X_0/Y_0, \mathcal{J}_{Y_0}) \quad \text{i.e. } dy = 0
\]
\uncertain{d'où} \ill{} \uncertain{cobord} \ill{} \ill{} $f$, et
\ill{} \ill{} que $p_1(f + \lambda) = p_2(f + \lambda)$, \ill{} \ill{} que
$p_2(\lambda) - p_1(\lambda) = \uncertain{-y}$, i.e. $d(\lambda) = y$.

\page{7}

\begin{tikzcd}
Y & X \arrow[l] & R \arrow[l, bend left=8] \arrow[l, bend right=8] \\
Y_0 \arrow[u] & X_0 \arrow[l, hook'] \arrow[u] & R_0 \arrow[l, bend left=8] \arrow[l, bend right=8] \arrow[u]
\end{tikzcd}

\note{la flèche verticale de gauche est repassée à l'encre foncée et son
sens n'est pas sûr\,; la verticale du milieu n'a pas de pointe visible}

D'ailleurs, \struck{pour $\mathcal{J}_{Y_0}$} \uncertain{$Y$} \ill{}
\uncertain{est} \ill{} \ill{}, $Y_0$ \ill{} \ill{} \ill{} \ill{}
\uncertain{défini} par un \ill{} \ill{} nilpotent \uncertain{isomorphe}\,:
$\mathcal{J}_{Y_0}$. \ill{} \ill{} \ill{} \ill{} \ill{}
\[
  \mathcal{J}_{Y_0} X = \mathcal{J}_{X_0} \simeq
  \mathcal{J}_{Y_0} \otimes_{\mathcal{O}_{Y_0}} \mathcal{O}_{X_0}
\]
\ill{} \ill{} $X$, i.e. $X_0 \to Y_0$ \ill{} \ill{}

\ill{} $X \to Y$ \uncertain{est} \ill{} \ill{} $Y$. \ill{} \ill{} \ill{}
\uncertain{que} $X_0/Y_0$ \ill{}, $X/Y$ \ill{} \emph{plat}. \ill{} $R \to X$
\ill{} plat sur $Y$, et $R \to X \times_Y X = P$ \ill{} \ill{}
\uncertain{isomorphisme}, \ill{} $R_0 \to P_0$ \ill{} \ill{}
\uncertain{isom.}, donc $R \to P$ \ill{} \ill{} isomorphisme, o.k.

\page{8}

Généralités sur la \uncertain{topologie fid. plate}

Les $\check{H}^p(S'/S, \mathcal{H}^q(F))$ pour $F$
\uncertain{quasi-cohérent}.

Applications aux relations d'équivalence ……
\note{feuillet brun, qui ne porte que ces trois lignes, en forme de plan\,;
il n'est pas certain qu'il se rapporte aux pages 1-7 qui le précèdent plutôt
qu'à celles qui suivent}

\section*{Prétopologies engendrées}
\note{titre de l'éditeur}

\page{9}

\textbf{Proposition.} $\mathcal{C}$ catégorie munie d'une prétopologie
$T$, \add{(avec produits fibrés)} $\mathcal{C}'$ sous-catégorie
\struck{\ill{}} pleine de $\mathcal{C}$ \struck{\ill{} \ill{}}
\struck{système de} \ill{} qui \uncertain{soit} \uncertain{stable} par $T$
\struck{\ill{} \ill{}} \struck{isomorphismes et produits fibrés}\,:

(a) pour tt $X \in \mathrm{Ob}\,\mathcal{C}$, existe famille couvrante
$U_\alpha \to X$ pour $T$, \ill{} les $U_\alpha \in \mathrm{Ob}\,\mathcal{C}'$\,],
\struck{\ill{}}

(b) $\mathcal{C}'$ stable par isomorphismes et produits fibrés, \struck{\ill{}}
\struck{\ill{} de \ill{} pour la}

$T'$ une \add{prétopologie} \struck{prétopologie} sur $\mathcal{C}'$\,!
\struck{\ill{} \ill{} \ill{} \ill{} \ill{}}

\struck{(c) stable par \ill{} \ill{} \ill{} \ill{}, \ill{} \ill{} \ill{} \ill{}
\ill{} \ill{} \ill{} \ill{} \ill{} identités \ill{} \ill{} \ill{} \ill{}
isomorphismes \ill{} \ill{} \ill{} $T'$. Soit $\mathcal{C}$ la topologie
engendrée par $T$, $T'$. Soit $X_i \to X$ une famille de flèches
\ill{} \ill{} \ill{} \ill{} couvrante pour $T$, \ill{} $f$. \ill{}}
\note{tout le bloc (c), avec le début de démonstration, est biffé de
plusieurs grandes croix et de traits ondulés\,; la lecture en est très
incertaine}

On \uncertain{va} que la situation
\[
\begin{array}{ccc}
 & & U_{ij\mu} \\
 & & \downarrow \\
U_i & \longleftarrow & U_{ij} \\
\downarrow & & \\
U & &
\end{array}
\]
où $U, U_i, U_j, U_{ij\mu} \in \mathrm{Ob}\,\mathcal{C}'$, $(U_i \to U)$
couvrante pour \struck{\ill{}} $T'$, et $(U_{ij\mu} \to U_{ij})$
\ill{} \ill{}, $(U_{ij}) \to U_i$ couvrante pour $T$, on \uncertain{veut}
\ill{} \ill{} \ill{} sous-famille de $U_{ij\mu} \to U$ couvrante pour $T'$.
\note{le $U_i$ du diagramme surcharge une autre lettre}

\page{10}

Soit $\mathcal{T}$ la topologie engendrée par $T$, $T'$, \ill{} \ill{} \ill{}
$X_i \to X$ une famille de morphismes de $\mathcal{C}$. Pour qu'elle soit
$\mathcal{T}$-couvrante, il faut et suffit qu'il existe une famille
$T$-couvrante $U_\alpha \to X$, \add{$U_\alpha \in \mathrm{Ob}\,\mathcal{C}'$,}
et pour tt $\alpha$ une famille $T'$-couvrante $U_{\alpha\beta} \to U_\alpha$,
et pour tt $\beta$ un indice $i$ et un $X$-morphisme
$U_{\alpha\beta} \to X_i$. \add{\uncertain{cor.}} Si
$X = \struck{U} \in \mathrm{Ob}\,\mathcal{C}'$, il faut et suffit aussi qu'il
existe $U_\alpha \to X$ couvrante pour $T'$, et pour tt $\alpha$ un $i$
\struck{\ill{}} un $X$-morphisme $U_\alpha \to X_i$.

\begin{tikzcd}[column sep=small, row sep=small]
X_i \arrow[d] & U_{\alpha\beta} \arrow[l] \arrow[d] \\
X & U_\alpha \arrow[l]
\end{tikzcd}

\marginal{N.B. la \uncertain{condition} \ill{} \ill{} la proposition, \ill{}
\ill{} (c)}\note{en marge gauche, en travers, séparé du texte par un trait}

\emph{Dém.} Si cette condition est remplie, alors $U_{\alpha\beta} \to X$ est
couvrante pour $\mathcal{T}$ par transitivité, et par suite $X_i \to X$ l'est
par saturation. Il reste à prouver que les familles du type envisagé
satisfont les conditions 1), 2), 3), 4). Pour 4) c'est trivial, pour 3)
aussi [prenons une famille couvrante \add{si $X' \rightrightarrows X$,}
\struck{$U_\alpha$} de $X$ \ill{} les $U_\alpha \in \mathrm{Ob}\,\mathcal{C}'$
--- existe par hypothèse (a) \struck{$\mathcal{C}'$ \ill{} de
$(\mathcal{C}, T)$}, $(U_{\alpha\beta})_\beta$ réduit à\,:
$U_\alpha \xrightarrow{\mathrm{id}} U_\alpha$, et $U_\alpha \to X'$ déduit de
$U_\alpha \to X$. Pour \uncertain{1)}\,\ill{}),

\page{11}

\struck{Alors} Pour tt $\alpha$, il existe par construction
\uncertain{$i = i(\alpha)$} \struck{tel} et un $X$-morphisme
\struck{$U_{\alpha\beta} \to$} $V_{\alpha\lambda} \to X_{\uncertain{i}}$,
ok.

\textbf{Cor. 2} Soit $F : \struck{\ill{}}\,\mathcal{C}^{\circ} \to
(\mathrm{Ens})$, \uncertain{pour que} $F$ soit un faisceau
\add{pour $\mathcal{T}$}, il faut et suffit que $F$ \struck{\ill{} \ill{}
\ill{}} soit un faisceau pour $T$, et \struck{\ill{}} induise sur
$\mathcal{C}'$ un faisceau pour $T'$\,!

\marginal{En effet, \ill{} \ill{} \ill{} \ill{}, \ill{} \ill{} \ill{} \ill{}
\ill{} \struck{\ill{}} \ill{} $F$ \ill{} \ill{} topologie, \ill{} pour $T$,
$T'$, \ill{} $\mathcal{T}$. OK}\note{en marge gauche, en travers, relié au
texte par une accolade}

\textbf{Proposition 2.} Considérons sur $(\mathrm{Sch}) = \mathcal{C}$ la
prétopologie $T$ de Zariski. Considérons la \struck{\ill{}} sous-catégorie
pleine $\mathcal{C}'$ de $\mathcal{C}$ formée des schémas affines, et la
prétopologie suivante\,: une famille de $U_i \to U$ est couvrante, si elle
est finie, \add{surjective et} \struck{\ill{}} plate [i.e. les $U_i$ sont en
nombre fini, leur somme recouvre $U$, et les $U_i \to U$ sont plats]. Soit
$\mathcal{T}$ la topologie \add{sur $\mathcal{C}$} engendrée par $T$, $T'$.

\marginal{Ex. Cas $(\mathrm{Sch})$, $T$\,: Zariski-topologie, $T'$\,:
$f : U \to V$ \ill{} \ill{}, $U$, $V$ affines}\note{entouré, en marge
gauche, en travers\,; au-dessous, toujours en travers\,: « \ill{} plate
\ill{}\,: plate de présentation finie, plate \ill{} \ill{}, \uncertain{étale}, »,
énumération de variantes possibles de $T'$}

(i) Soit $X_i \to X$ une famille de morphismes dans $\mathcal{C}$. Pour
qu'elle soit couvrante, il f. et s. que $X$ puisse se recouvrir par des
ouverts affines $U_\alpha$ et que pour tt $\alpha$ $\exists$ une famille finie
plate $U_{\alpha\beta} \to U_\alpha$, les $U_{\alpha\beta}$ affines,
\struck{\ill{}} et pour tt $\alpha, \beta$ un $i = i(\alpha, \beta)$ et un
$X$-morphisme $U_{\alpha\beta} \to X_i$.

\begin{tikzcd}[column sep=small, row sep=small]
X_i \arrow[d] & U_{\alpha\beta} \arrow[l] \arrow[d] \\
X & U_\alpha \arrow[l]
\end{tikzcd}

\page{12}

(iii) Un morphisme fid. plat quasi-compact $f : X' \to \struck{S}\,X$ est
couvrant pour $\mathcal{T}$. De même un \add{ouvert, p.ex. un} morphisme fid.
plat \struck{loc. de présentation finie}, morphisme fid. plat loc. de prés.
finie\,; plus généralement, si on a une famille de morphismes $X_i \to X$
\emph{ouverts}, \emph{plats}, dont les images recouvrent $X$, cette famille
recouvre $X$ pour $\mathcal{T}$.
\note{le paragraphe (iii) est rattaché par un trait de marge au (ii) qui
suit}

(ii) Soit $F : (\mathrm{Sch})^{\circ} \to (\mathrm{Ens})$ un foncteur. Pour
que $F$ soit un faisceau pour $\mathcal{T}$, il f. et s. qu'il satisfasse les
conditions suivantes\,:

a) $F$ est un faisceau pour la topologie de Zariski,
i.e. « $F$ de nature locale »

b) Pour tt $S' \to S$ fid. plat, avec $S$ et $S'$ affines, on a exactitude
dans
\[
  F(S) \to F(S') \rightrightarrows F(S' \times_S S')
\]

\emph{Dém.} (i) résulte de Prop.~1.

(ii) \uncertain{Résulte trivialement}, par 2 \ill{}, Cor.~1 corollaire 2,
\ill{}\,; \struck{\ill{} \ill{} \ill{} \ill{} \ill{} par $F$ \ill{} \ill{}}
\struck{pour les morphismes fid. plats et quasi-compacts} [idem
\uncertain{données} \ill{} \ill{} sont stables par
\note{les deux dernières lignes sont barrées d'un trait ondulé et
encadrées ; la phrase s'arrête au bas du feuillet}

\page{13}

\note{suite de la démonstration de la page~10 (vérification des
axiomes)}

\begin{tikzcd}[column sep=small, row sep=small]
X_i \arrow[d] & U_{\alpha\beta} \arrow[l] \arrow[d] & \\
X & U_\alpha \arrow[l] & \\
X'_i \arrow[uu, bend left=20] \arrow[d] & U'_{\alpha\beta} \arrow[l] \arrow[d] & U'_{\alpha\beta\lambda} \arrow[l] \arrow[d] \\
X' \arrow[uu, bend left=30] & U'_\alpha \arrow[l] \arrow[uu, bend right=20] & U'_{\alpha\lambda} \arrow[l]
\end{tikzcd}

\note{le dessin superpose les deux carrés, celui de $X'$ décalé vers le bas
à droite, avec des flèches obliques de $X'_i$, $X'$, $U'_\alpha$ vers
$X_i$, $X$, $U_\alpha$ ; on l'a déplié. $X'_i$ est sans doute
$X_i \times_X X'$}

\ill{} \ill{} que les $U'_\alpha$ recouvrent $X'$ pour $T$\,; \struck{\ill{}}
pour tt $\alpha$, \add{recouvrons \uncertain{pour $T$}} $U'_\alpha$ par des
$U'_{\alpha\lambda} \in \mathrm{Ob}\,\mathcal{C}'$, \uncertain{possible} par
hyp.~(a), d'où des $U'_{\alpha\beta\lambda}$, pour $\alpha$ et $\lambda$
fixés, les $U'_{\alpha\beta\lambda}$ recouvrent \add{\uncertain{puisque}
(b), en déduisant de $U_{\alpha\beta} \to U_\alpha$ par changement de base,}
$U'_{\alpha\beta}$ pour $T'$\,; \add{\uncertain{de plus} pour $\alpha$,
$\lambda$ variables,} les $U'_{\alpha\lambda}$ recouvrent $X'$ pour $T$, par
transitivité. OK.

Axiome 1 de transitivité, la construction \ill{}. Soient\,: l'hyp.\
\add{\uncertain{sur}} $X_i \to X$ \struck{\ill{}} des $U_\alpha$,
$U_{\alpha\beta}$, et pour tt $\alpha, \beta$ --- $U_{\alpha\beta} \to X_i$,
\add{$i = i(\alpha, \beta)$,} \ill{} \ill{} \ill{} \ill{} \ill{} \ill{}, la
famille des $X_{ij} \times_{X_i} U_{\alpha\beta}$ (\uncertain{qu'elle})
satisfait aux conditions envisagées, donc $\exists$
$U_{\alpha\beta\gamma} \to U_{\alpha\beta}$ couvrant pour $T$, les
$U_{\alpha\beta\gamma} \in \mathrm{Ob}\,\mathcal{C}'$, et \struck{pour} tt
$\gamma$ une famille $U_{\alpha\beta\gamma\delta} \to U_{\alpha\beta\gamma}$
couvrante pour $T'$, et un $j = j(\gamma, \delta)$ et $U_{\alpha\beta}$-morphisme
$U_{\alpha\beta\gamma\delta} \to \struck{X}\,X_{ij} \times_{X_i}
U_{\alpha\beta}$. Par (c), \ill{} \ill{} pour tt $\alpha$, \uncertain{extraire}
de $U_{\alpha\beta\gamma} \to U_\alpha$ une famille couvrante par $T'$, soit
$V_{\alpha\lambda} \to U_\alpha$.

\begin{tikzcd}[column sep=small, row sep=small]
X_{ij} \arrow[d] & X_{ij} \times_{X_i} U_{\alpha\beta} \arrow[l] \arrow[d] & U_{\alpha\beta\gamma\delta} \arrow[l] \arrow[d] \\
X_i \arrow[d] & U_{\alpha\beta} \arrow[l] \arrow[d] & U_{\alpha\beta\gamma} \arrow[l] \\
X & U_\alpha \arrow[l] &
\end{tikzcd}

\note{$U_\alpha$ est entouré sur le dessin. La condition (c) invoquée ici
est celle qui est biffée à la page~9}

\page{14}

\struck{\ill{} engendrée en fait dans $\mathcal{C}$, et engendre
$\mathcal{T}$ \ill{} la prétopologie de Zariski). \ill{} \ill{} ceci, il
suffit de}
\note{deux lignes barrées d'un trait ondulé, suite de la phrase interrompue
au bas de la page~12}

\begin{tikzcd}[column sep=small, row sep=small]
S' \arrow[d] & \\
S & S_\alpha \arrow[l] \arrow[lu, bend right=30, no head]
\end{tikzcd}

\ill{} qui\,: condition \uncertain{a)}, $F$ \struck{\ill{}} transforme
\uncertain{sommes} \ill{} en produits, ce qui permet d'appliquer la
condition \uncertain{b)} du corollaire 1, \uncertain{au} produit de
\uncertain{deux} \uncertain{recouvrements} \ill{} \ill{} \ill{} deux
morphismes.

\struck{\ill{}} \add{(iii)} \uncertain{par des} \uncertain{produits}……
\note{le reste du feuillet est blanc}

\section*{Descente fidèlement plate}
\note{inscrit à l'encre, en haut à droite et souligné, sur le feuillet brun
de la page 15, qui ne porte rien d'autre et ne reçoit pas de numéro de page ;
le feuillet 16 est blanc}

\page{17}

\struck{1} A) Application à la comparaison entre les variétés
\uncertain{étales} holomorphes

B) Application à \struck{la construction} \add{la \uncertain{factorisation}
canonique} \uncertain{par}

\textbf{Proposition.} Soit $f : X \to Y$ un morphisme
\add{\uncertain{univ. ouvert}} fini \struck{\ill{}} avec $Y$ loc.
noethérien. \struck{Pour} $y \in Y$, soit $d(y)$ le \uncertain{rang} de
\uncertain{sa} \uncertain{géométrique} de $f^{-1}(y)$ (= degré
séparable de $f^{-1}(y)$ sur $k(y)$)]. \emph{Supposons} $d(y)$ constant.
Alors $f$ se factorise en
\[
  X \xrightarrow{f_1} X^{\ast}_0 \xrightarrow{f_2} Y
\]
où $f_1$ est \add{fini} \uncertain{surjectif} radiciel, et $f_2$ est un
morphisme \uncertain{fini} \emph{étale}. \add{$f_2$ Alors} \struck{\ill{}}
\uncertain{Tout} $Y$-morphisme $X \to Z$, \uncertain{avec} $Z$ étale sur $Y$,
se factorise de façon unique en $X \to X' \to Z$.

On \uncertain{peut} \uncertain{supprimer} \add{$Y$ connexe \uncertain{dans}}
$d(y)$ constant, \uncertain{sauf} \ill{}. [\,\uncertain{Quitte} à
\uncertain{remplacer} $Y$ par une partie \ill{} \uncertain{finie} \ill{} et
\uncertain{fermée} \ill{}, \ill{} tt \ill{}\,]
\note{la fin de la page, entre crochets, est très rapide ; l'exposant de
$X_0$ dans la factorisation est lu $\ast$, la lettre n'est pas sûre}

\page{18}

\note{la page entière est biffée de deux longues diagonales ; elle se lit
néanmoins}

Soit $I = [1, \ldots, n]$. Considérons
\[
  Y' = \mathrm{Hom}(I, X) = X^{I}
\]
On a donc un $Y$-morphisme \uncertain{canonique}
\[
  Y' \times I \to X \qquad \text{i.e. un } Y'\text{-morphisme
  \uncertain{canonique}}
\]
\[
  f : \struck{\ill{}}\,Y' \times I \to X'
\]
[correspondant aux $n$ projections $\mathrm{pr}_i : Y' = X^I \to X$]
\note{$X'$ désigne sans doute $X \times_Y Y'$}

Pour tout $y' \in Y'$, on en conclut une \uncertain{application}
\[
  I \times \mathrm{Spec}\,k(y') \to X' \times_{Y'} \mathrm{Spec}\,k(y')
\]

\textbf{Lemme.} Soit \add{$Y''$} \struck{\ill{}} \uncertain{l'ensemble} des
pts de $Y'$ où le morphisme précédent est bijectif \add{\uncertain{ou}}\,:
la \uncertain{fonction} \ill{} \ill{} \uncertain{ferme}, et
\struck{il s'applique} \add{$Y'' \to Y$} \struck{bijectivement} surjectif
\struck{\ill{} \ill{} \ill{} \ill{}} \add{$Y'' \to$ \uncertain{connexe}.}

\struck{L'un des pts} \uncertain{L'un} des pts de $I \times \mathrm{Spec}\,k(y)
\to X'_y$ est géométriquement \emph{injectif}, i.e. \ill{} \ill{}
\uncertain{est} \emph{radiciel}, \uncertain{ou} \ill{} \ill{} \ill{}
un \uncertain{monomorphisme} [= \uncertain{immersion fermée}], \ill{}
\uncertain{\ill{}} (\uncertain{décrit}) \uncertain{que} \ill{} \ill{}
\ill{}.
\marginal{$\struck{Y'}\,Y' \times I$ \uncertain{est} plat \struck{propre}
sur $Y'$}\note{entouré, en interligne au-dessus de la dernière ligne}
\marginal{\uncertain{décrit} \uncertain{explicitement} \struck{est} de
2 \ill{}}\note{en marge gauche, en travers}

\page{19}

Si $X \xrightarrow{f} Y$ admet une section $s$, \struck{\ill{}} alors $s(Y)$
est à la fois \struck{\ill{}} fermé et \uncertain{ouvert}. En effet,
\struck{on voit de suite que c'est une} \uncertain{réunion} de
\uncertain{composantes} irréductibles de $X$ [$X \to Y$ étant fini]\,; il
\uncertain{faut} \uncertain{montrer} que si $X''$ est \struck{la} \ill{}
\uncertain{réunion} des autres composantes irréductibles, alors
$X' \cap X'' = \emptyset$. Or \struck{\ill{} \ill{}}
\struck{l'image \ill{} \ill{} \ill{}} \ill{} de $Y$, \ill{}
\add{les composantes irréductibles} \struck{$d''(y) = n - 1$, donc}
\note{le passage de « il faut montrer » à « donc » est encadré, et barré
de deux traits obliques ; $X'$ désigne sans doute $s(Y)$}

\marginal{$X$ / $Y$}\note{en marge gauche, isolés}

\ill{} \uncertain{Pour} \uncertain{un} \uncertain{domaine} \ill{} \ill{}
\uncertain{supposer} \ill{} \ill{} $Y$ \uncertain{est} irréductible
[car si les $Y_i$ sont les composantes de $Y$, il suffit de montrer que
$s(Y_i)$ \uncertain{est ouvert} \ill{} \uncertain{dans} $f^{-1}(Y_i)$]

\begin{tikzcd}[column sep=small, row sep=small]
X \arrow[d, bend left=20] & f^{-1}(Y_i) \arrow[d, bend left=20] \\
Y \arrow[u, bend left=20, "s"] & Y_i \arrow[u, bend left=20, "s_i"]
\end{tikzcd}

\note{croquis de marge, redessiné : $X = \bigcup_i f^{-1}(Y_i)$ au-dessus
de $Y = \bigcup_i Y_i$, avec $s$ et $s_i$ en montant}

…\ et alors les composantes irréductibles de $X$ sont
\emph{disjointes} [\uncertain{vérif}], \struck{\ill{}} \ill{} \ill{}
la conclusion.

Donc \ill{} \ill{} $X = X_{(1)} \amalg X^{(1)}$, où
\[
  d(X_{(1)}/Y) = 1, \qquad d(X^{(1)}/Y) = n - 1 .
\]

\page{20}

\note{suite, semble-t-il, du lemme de la page~18}

et que $y' \in Y'' \Leftrightarrow$ il existe un voisinage $U$ de $y'$ dans
$Y'$ tel que $f : U \times I \to X'|U$ soit une \struck{\ill{}} immersion
fermée. Le complémentaire de $Y''$ est formé des pts tels que
\[
  I \times \mathrm{Spec}\,k(y') \to X' \times_{Y'} \mathrm{Spec}\,k(y')
\]
soit géométriquement non injectif,
\note{la phrase s'arrête là ; le reste du feuillet est blanc}

\end{document}
